
\chapter{Analisi schematica dei requisiti}

\section{Obiettivo del sistema}
Realizzare un'applicazione client/server per la gestione di una chat testuale, denominata Ruggine, che consenta la comunicazione tra utenti e la creazione di gruppi per la condivisione di messaggi.

\section{Requisiti funzionali}
\begin{itemize}
	\item \textbf{Iscrizione utente}: Al primo avvio, l'utente deve iscriversi tramite una richiesta al server, ottenendo così le credenziali necessarie per accedere al sistema.
	\item \textbf{Login}: L'autenticazione consente all'utente di accedere alle funzionalità della chat, garantendo la sicurezza delle informazioni.
	\item \textbf{Scambio di messaggi}: Il sistema permette l'invio e la ricezione di messaggi testuali in tempo reale tra utenti e gruppi.
	\item \textbf{Creazione gruppi}: Ogni utente ha la possibilità di creare nuovi gruppi di chat, favorendo la comunicazione tematica.
	\item \textbf{Gestione inviti}: L'accesso ai gruppi avviene solo tramite invito; il sistema gestisce l'invio, la ricezione e l'accettazione/rifiuto degli inviti.
	\item \textbf{Compatibilità multipiattaforma}: L'applicazione deve funzionare su almeno due tra Windows, Linux, MacOS, Android, ChromeOS, iOS, garantendo ampia accessibilità.
	\item \textbf{Logging}: Il server produce periodicamente (ogni 2 minuti) dati di log riguardanti la CPU utilizzata, utile per il monitoraggio delle prestazioni.
	\item \textbf{Report dimensione eseguibile}: Nel report finale deve essere inclusa la dimensione del file eseguibile.
\end{itemize}

\section{Requisiti non funzionali}
\begin{itemize}
	\item \textbf{Prestazioni}: Il sistema è progettato per ridurre il consumo di CPU e la dimensione dell'applicativo, garantendo efficienza anche su dispositivi con risorse limitate.
	\item \textbf{Sicurezza}: L'autenticazione degli utenti e la protezione dei dati sono gestite con tecniche robuste, per prevenire accessi non autorizzati e garantire la privacy.
	\item \textbf{Scalabilità}: L'architettura consente di gestire efficacemente un numero crescente di utenti e gruppi, senza degradare le prestazioni.
	\item \textbf{Usabilità}: L'interfaccia utente è pensata per essere semplice e intuitiva, facilitando l'accesso alle funzionalità principali.
	\item \textbf{Portabilità}: Il codice è strutturato per essere facilmente adattabile alle diverse piattaforme, favorendo la diffusione dell'applicazione.
	\item \textbf{Manutenibilità}: La progettazione modulare facilita l'estensione e la manutenzione futura del sistema.
\end{itemize}

\section{Vincoli}
\begin{itemize}
	\item \textbf{Piattaforme supportate}: Il sistema deve essere compatibile con almeno due tra Windows, Linux, MacOS, Android, ChromeOS, iOS, come richiesto dalla traccia.
	\item \textbf{Dimensione eseguibile}: La dimensione del file eseguibile deve essere esplicitamente riportata nel documento finale.
	\item \textbf{Log CPU}: Il server deve generare un log con dettagli sul tempo di CPU ogni 2 minuti, per consentire il monitoraggio delle risorse.
\end{itemize}

\section{Casi d'uso principali}
\begin{itemize}
	\item \textbf{Iscrizione utente}: L'utente si registra al sistema, ottenendo le credenziali per accedere alle funzionalità.
	\item \textbf{Login/logout}: L'utente effettua l'accesso per utilizzare la chat e può uscire in qualsiasi momento, terminando la sessione.
	\item \textbf{Messaggistica}: Gli utenti possono inviare e ricevere messaggi in tempo reale, sia in conversazioni private che di gruppo.
	\item \textbf{Creazione gruppo}: È possibile creare nuovi gruppi di chat, invitando altri utenti a partecipare.
	\item \textbf{Gestione inviti}: Il sistema consente di inviare, ricevere e gestire inviti ai gruppi, con possibilità di accettare o rifiutare.
	\item \textbf{Partecipazione a gruppi}: Una volta accettato l'invito, l'utente può interagire nei gruppi, scambiando messaggi con i membri.
	\item \textbf{Visualizzazione log}: Gli utenti e gli amministratori possono consultare i log e le statistiche di sistema per monitorare l'attività.
\end{itemize}